\section{Discussion}
\label{sec:discussion}
\subsection{What limits the magnification of a lensed GW image?}
Our results separate five sources of magnification error. (i) \emph{Lens-model freedom at fixed image configuration.} This dominates for realistic lenses: from SIS to SIE, SIE+shear and EPL+shear the geometry-only error grows from 3 to 7, 12 and 20--21 per cent for ET. The radial density slope is the most important unobserved parameter, as expected from the slope--magnification degeneracy familiar from time-delay cosmography. (ii) \emph{Proximity to a caustic.} Errors grow roughly in proportion to $\mu$, so the most strongly magnified---and, for current detectors, most likely to be detected---images are the hardest to correct. (iii) \emph{The GW apparent distance.} Under the catalogue-calibrated error model, the inclination--distance degeneracy limits each image's apparent distance to $\sigma_{\ln D}\approx0.22$, and this dominates the error of the corrected luminosity distance. Joint PE of both images with higher harmonics gives apparent-distance errors smaller by a factor 2.3 for ET and 1.5 for O5 (Section~\ref{sec:jointpe}), so this conclusion depends on the GW analysis and should not be read as a fundamental floor. (iv) \emph{A wrong lens prior.} An estimator trained on an over-simplified lens population produces confidently wrong intervals. This is the most dangerous failure mode, because nothing in the data flags it. (v) \emph{An implicit distance scale.} An estimator that receives the delay or the SNRs together with redshifts gains accuracy by learning the simulated cosmology, and its output then shifts with the assumed $H_0$ (Table~\ref{tab:response}).

These findings suggest a practical strategy for the first lensed BBH with an identified host. Image geometry and the GW relative magnification should be combined with a lens prior that is at least as flexible as an EPL plus external shear with an empirically calibrated slope distribution. Intervals should be calibrated on that flexible prior, because calibration groups defined by observables do not reproduce per-family calibration (Section~\ref{sec:calibration}). Inputs that carry a distance scale should be used only when that scale is part of the analysis rather than an unexamined assumption. Any additional handle on the radial slope---stellar kinematics with a dynamical model, extended-arc modelling of the host, or the GW time delays themselves combined with an independent time-delay distance---should be added, because it targets the dominant error term.

\subsection{Mass-sheet degeneracy and cosmology dependence}
Every result above is conditional on lens populations without an external mass sheet. For estimators that use only scale-free inputs, a sheet with $\lambda\neq1$ multiplies the estimate relative to the truth by exactly $\lambda^2$, which we verified numerically (Table~\ref{tab:response}). For a typical external convergence $|\kappa_{\rm ext}|\lesssim0.05$, this bias of up to 10 per cent in $\mu$ (5 per cent in $d_L$) is comparable to our SIE errors and below those of the EPL family. It would, however, be coherent across a population of lensed events if environments were systematically over- or under-dense. Estimators that use the delay respond to $\lambda$, but only partially (Table~\ref{tab:response}), so they do not remove the bias either. The Bayesian reconstruction makes the trade-off explicit: with $\kappa_{\rm ext}$ free, the scale-free information leaves it unconstrained and the magnification posterior widens accordingly, whereas the delay together with an assumed cosmology constrains $(1-\kappa_{\rm ext})D_{\Delta t}$ (Section~\ref{sec:gwem}). Line-of-sight convergence estimates from galaxy counts or ray tracing, as in time-delay cosmography, can be folded into the prior. Any constraint that uses an assumed cosmology, however, makes the corrected distance circular for cosmography, so the scale-free results are the relevant benchmark for that application.

\subsection{Limitations}
We assumed that the GW images have been associated with each other and with the EM images of an identified host, and that the image positions of the BBH within the host are reconstructed to the stated astrometric precision. The BBH itself is not observed in the EM data, so these positions are the result of an upstream reconstruction of the extended host and the GW information. Its uncertainty may be larger, correlated between images and multimodal, and it must not reuse the GW measurements that enter the magnification inference. Section~\ref{sec:gwem} shows how the magnification degrades with poorer astrometry. A pixel-level reconstruction of lensed hosts with off-centre BBHs, propagated into the GW-source image positions, is needed to replace this assumption; we leave it to future work. The association steps are probabilistic, and the magnification posterior must marginalise over candidate associations and lens reconstructions \citep{hannuksela2020localizing,wempe2022lensing,uronen2026galacticcentres}. The benchmark itself uses the $(2,2)$ mode of an aligned-spin waveform, an SNR-based amplitude model and analytic scalings for the sky and orientation errors. The O4a injections, which recover each image with the exact template, test the relative-amplitude noise prescription in the LIGO detectors only; the joint PE of Section~\ref{sec:jointpe} tests the SNR-based measurement model with higher harmonics and free source parameters, but in Gaussian noise and with aligned spins. Lens populations do not include satellites, line-of-sight structures or redshift evolution of lens properties, and wave-optics effects are negligible for the galaxy-scale lenses considered here but would matter for lower-mass lenses. Each of these simplifications should be relaxed before the method is applied to data.

\section{Conclusions}
\label{sec:conclusions}
We have benchmarked the accuracy and reliability of absolute-magnification inference for strongly lensed GW images from GW and EM observables. The benchmark uses GWTC-4.0 source populations, four galaxy-lens families, O5-like and ET networks, GW distance errors calibrated on 167 GWTC events, and tests of the GW measurement model in public O4a noise and by joint parameter estimation. Our main conclusions are as follows.
\begin{enumerate}
\item Image geometry measured relative to the lens-light centre is the most informative observable. Combined with the GW relative magnification and arrival order---inputs that are invariant under the mass-sheet transformation and contain no distance scale---it determines the magnification of an ET image to a median 3.0, 6.9, 10.2 and 18.8 per cent for SIS, SIE, SIE+shear and EPL+shear lenses (HistGB trained on the correct family; NPE trained on all families gives 3.3--19.6 per cent). For the more magnification-biased O5 sample the errors are 10.8--25.5 per cent.
\item The dimensional time delay and the SNRs add little on their own. With redshifts they reduce the errors further, to 1.7--13.6 per cent for NPE, but only by learning the distance scale of the training cosmology: the inferred magnification then shifts by 0.7 per cent per per cent change in $H_0$. Such estimates are not appropriate for cosmography.
\item GW data alone, or the apparent distance combined with a known source redshift, determine the magnification much less precisely (21--26 and about 41 per cent for realistic lenses), the latter because of the catalogue-calibrated inclination--distance degeneracy.
\item The Gaussian relative-amplitude model holds in real O4a noise for exact-template recovery (standardised scatter 1.02), apart from one glitch-contaminated pair among 321. Joint parameter estimation of 120 image pairs from simulated detector data, with higher harmonics and free source parameters, supports the SNR-based measurement model: with the host sky position the 90 per cent intervals of the relative magnification and the apparent distance contain the truth in 83--92 per cent of systems, and the PE uncertainties are smaller than those of the model. With a free sky position, ET posteriors from a single site have several separated solutions that the sampler does not always find, so the host position is needed in practice.
\item Within a lens family, the calibrated (conformal) 90 per cent intervals contain the truth 90 per cent of the time. An estimator trained on all families needs the family label to be calibrated per family; groups of events defined by observables do not substitute for it. Calibrating on the most flexible (EPL+shear) population gives at least 90 per cent coverage in every family at 40 per cent wider intervals. Near caustics, the intervals contain the truth too rarely for the truly highly magnified images (44--57 per cent for $\mu_A>20$), and calibrating in bins of predicted magnification does not fix this.
\item The lens prior is the dominant systematic. SIS-trained models are confidently wrong for realistic lenses (coverage 20--27 per cent), whereas a flexible EPL+shear prior transfers conservatively to the simpler families we tested. Omitting radial-slope freedom drops interval coverage to 55 per cent.
\item With the catalogue-calibrated apparent-distance errors, the GW apparent distance dominates the error of the corrected luminosity distance for typical lensed BBHs (22--24 per cent total for ET, of which 1.5--9.5 per cent comes from the magnification). Joint parameter estimation with higher harmonics reduces the apparent-distance error by a factor 2.3 for ET (1.5 for O5), so this balance depends on the GW analysis.
\item A conventional Bayesian GW+EM reconstruction with the same information reaches 8.1 (ET) and 13.7 (O5) per cent, comparable to the NPE (10.5 and 13.8 per cent), with 40--45 per cent narrower intervals that contain the truth slightly too rarely for ET. Summarising the GW posterior by a Gaussian in the relative magnification reproduces the full analysis of the detector data (median shift $0.15\sigma$ or less, equal widths) when the GW posterior is converged. The result degrades quickly with the astrometry of the GW-source images, doubling at 0.1\,arcsec.
\end{enumerate}
The code, configuration and real-data calibration are public. The results are conditional on reconstructed GW-source image positions and on the absence of an external mass sheet; pixel-level reconstruction of lensed hosts with off-centre BBHs and line-of-sight convergence are the most important next steps.
